\section{Introductions}

Suppose $M$ is an $n$-dimensional K\"ahler manifold and $L$ is an ample
line bundle over $M$. Let the K\"ahler form of $M$ be $\omega_g$ and the
Hermitian metric of $L$ be $H$. We assume that $\omega_g$ is the curvature
of $H$, that is, $\omega_g=Ric(H)$. 
The K\"ahler metric of
$\omega_g$ is called a polarized 
\ka metric on $M$.

Using $H$ and $\omega_g$,
for any positive integer $m$,
 $H^0(M,L^m)$ becomes a Hermitian inner product space. 
We use the following notations: suppose that $S,T\in H^0(M,L^m)$. Let
$<S,T>_{H^m}$ be the pointwise inner product and
\[
(S,T)=\int_M<S,T>_{H^m}\frac{\omega^n_g}{n!}
\]
be the inner product of $H^0(M,L^m)$. Let
\[
||S||=\sqrt{<S,S>_{H^m}
}
\]
be the pointwise norm. In particular, $||S||(x)$ denotes the pointwise
norm at $x\in M$. Let
\[
||S||_{L^2}=\sqrt{(S,S)}
\]
be the $L^2$-norm of $S$. 

Let $\{S_1,\cdots,S_{d(m)}\}$ be an orthonormal basis of $H^0(M,L^m)$.
The quantity(see~\cite{T5})
\begin{equation}\label{quan}
\sum_{i=1}^{d(m)}||S_i||^2
\end{equation}
plays an important rule in K\"ahler-Einstein geometry and stability of
complex manifolds. 
If $m$ is sufficiently large, then there is a natural embedding
$\phi_m: M\rightarrow CP^{d(m)-1}$ by
$x\mapsto [S_1(x),\cdots,S_{d(m)}(x)]$. 
This follows from the Kodaira's embedding theorem.
The metric
$\frac 1m\phi^*\omega_{FS}$ on $M$ is called the Bergman metric. 
The following formula is an easy but key observation made by
Tian~\cite{T5}
in his proof of the covergence of the Bergman metrics:
\[
\frac 1m\pa\bar\pa\log\sum_{i=1}^{d(m)}||S_i||^2=\frac 1m
\phi^*\omega_{FS}-\omega_g.
\]


A lot of work has been done by several authors 
~\cite{T5},
~\cite{ru},~\cite{sz}, ~\cite{cat} and~\cite{Lu10}
on the estimates
of~\eqref{quan}. However, these works are concentrated on a
single manifold. On the other hand, in studying the stability and
K\"ahler-Einstein
geometry
of  manifolds, we need to study the behavior of~\eqref{quan} for
a
family of manifolds. Tian's work~\cite{T6} on the Calabi Conjecture shows
the
importance and non-triviality of the problem of giving lower bound
estimate 
of~\eqref{quan}
for a family of complex surfaces. 
In the  $n$-dimensional case, Tian~\cite{T9} proved that
~\eqref{quan} has
a positive lower bound depends on the
dimension $n$, the upper bound of the Betti numbers, the positve upper and
lower bound of the Ricci curvature and the $L^n$ norm of the Riemannian
sectional curvature.
In general, 
~\eqref{quan}  maybe a useful tool in studying algebraic fiberations over
a compact K\"ahler manifold.
The paper of interest to us are~\cite{MR83h:14025}
and~\cite{MR86j:14032}.


In this paper, we shall study the behavior of~\eqref{quan} on
Riemann surfaces. Even in the case of Riemann surfaces, the problem
of finding a uniform lower bound of~\eqref{quan} is nontrivial.
In fact, by
the counterexample in \S 3, we know that there is no uniform lower
bound in general.  We proved a partial uniform estimate in \S 4
which, I believe, is the ``right'' one in the sense that it gives all the
information on  stability of the Riemann surfaces. 


We also consider the coordinate ring of Riemann surfaces. For smooth
Riemann surfaces $M$ of genus $g\geq 2$, it is well  known that its
coordinate ring is
finitely
generated. That is, there is a positive integer $m_0$, such that
for any $S\in H^0(M, K_M^m)$ with $m>m_0$, we can find $U_i\in H^0(M,
K_M^{m_0}), i\in
I$ and $T_i\in H^0(M, K_M^{m-m_0}), i\in 
I$
such that
\[
S=\sum_{i\in I} U_iT_i,
\]
where $I$ is a finite set. In studying the behavior of 
Riemann surfaces near the boundary of the Teichm\"uller space, we need
some uniform estimates. In this paper, we give a uniform estimate
which will give us the information on  singular Riemann surfaces.  

The above setting is similar to that in the corona problem
in complex analysis.
The corona problem on the unit disk was studied by Carleson
in ~\cite
{MR25:5186}. Carleson's result
stimulates many ideas which proved to be useful for other problems.
An extensive discussion of the Carleson's corona theorem can be found
in~\cite{JBG}. 

Modifying Wolff's~\cite{WT} proof of Carleson's
theorem together with the $\bar\pa$-estimate, we give
a uniform corona estimate in the last section of this paper
as an application of Theorem~\ref{third}.
In order to obtain the result, we  take
 special care
to the points where the injective radius are small. 















The organization of the paper is as follows: in \S 2, we give a lower
bound of~\eqref{quan} in terms of the genus $g$ and the injective radius
$\delta$ of $M$. In \S 3, we give a counterexample which shows that the
lower bound must depend on $\delta$. In \S 4, we give the partial uniform
estimate. That is, a lower bound of ~\eqref{quan} at $x\in M$ depending
only on the injective radius $\delta_x$ of $x$. In \S 5, we solve the
uniform corona problem by the partial uniform estimate.






The main results of this paper are the following:

\begin{theorem}\label{main}
Let $M$ be a Riemann surface of the genus $g\geq 2$. Let $K_M$ be the
canonical line  bundle of $M$ endowed with a Hermitian metric $H$. Let 
the curvature $\omega_g$ of $H$ be positive. $\omega_g$ gives a \ka
metric of $M$. Let the curvature $K$ of $\omega_g$ satisfy
\[
-K_1\leq K\leq K_2
\]
for nonnegative constants $K_1, K_2\geq 0$ and let $\delta'$ be the 
injective radius of $M$. Let
\[
\delta=\min(\delta',\frac{1}{\sqrt{K_1+K_2}}).
\]
Then there is an absolute constant $C>0$ such that for $m\geq 2$,
\[
\sum_{i=1}^{d(m)}||S_i||^2\geq  e^{-\frac{Cg^3}{\delta^6}},
\]
where $\{S_1,\cdots,S_{d(m)}\}$ is an orthnormal basis of $H^0(M,K_M^m)$.
\end{theorem}


\begin{theorem}\label{second}
For any $\eps>0$ and $m\geq 2$, there is a Riemann surface
$M$
 of genus
$g\geq 2$ with the constant Gauss curvature $(-1)$ such that
\[
\inf_{x\in M}\sum_{i=1}^{d(m)}||S_i||^2\leq \eps.
\]
\end{theorem}


The theorem disproves the conjecture that
the absolute lower bound exists.

When the injective radius of $M$ goes to zero, the first eigenvalue 
and the Sobolev constant will also go to zero. In this case, Theorem
~\ref{main} gives no information. In the following theorem, we proved that
~\eqref{quan} has a lower bound which depends only on the local
information and is independent to the injective radius of $M$. 
For this reason, we call the result partial uniform estimate. 

\begin{theorem}\label{third}
Let $M$ be a Riemann surface of genus $g\geq 2$ and constant curvature
$(-1)$. Then there are absolute constants $m_0>0$ and $D>0$ such that for
any $m>m_0$
and any $x_0\in M$, there is a section $S\in H^0(M,K_M^m)$ with
$||S||_{L^2}=1$ such that
\begin{equation}
||S||(x_0)\geq\frac{\sqrt m}{D(1+\frac{1}{\sqrt m\delta_{x_0}^2}
e^{\frac{\pi}{\delta_{x_0}}})},
\end{equation}
where $\delta_{x_0}$ is the injective radius of $\delta_{x_0}$.
\end{theorem}


On the coordinate ring $\oplus_{j=0}^{\infty}
H^0(M, K_M^j)$, we have the following

\begin{theorem}
Let $M$ be a Riemann surface
as above.
Then there is an $m_0>0$ such that for any $m>m_0$ and $S\in
H^0(M,K_M^m)$,
there is a decomposition
\[
S=\sum_{i=1}^d S_i
\]
of $S_i\in H^0(M,K_M^m) (i=1,\cdots, d)$
such that
\begin{align}
\begin{split}
&||S_i||_{L^2}\leq C(m,m_0,g)||S||_{L^2}\\
&||S_i||_{L^\infty}\leq C(m,m_0,g)||S||_{L^\infty}
\end{split}
\end{align}
for $i=1,\cdots d$,
and
\[
S_i=T_iU_i
\]
for a basis $U_1,\cdots,U_d$ of $H^0(M,K_M^{m_0})$
and $T_1,\cdots,T_d\in H^0(M,K_M^{m-m_0})$.
\end{theorem}



{\bf Acknowledgment.} The author thanks Professor Tian for
suggesting me this problem and
 his help during 
the preparing of this paper. He also thanks Professor Phong for the
support
and some suggestions of the work.
